\documentclass[12pt]{article}

\usepackage[a4paper, total={6in, 8in}]{geometry}
\usepackage{textcomp}

\begin{document}
\setlength{\parindent}{0pt}

\def \t {\textrightarrow}
\def \v {\vspace{1em}}
\def \bi {\begin{itemize}}
\def \ei {\end{itemize}}
\def \s[#1] {\section*{#1}}
\def \ss[#1] {\subsection*{#1}}
\def \sss[#1] {\subsubsection*{#1}}

\s[Negli altri paesi]
\ss[Germania]
Hitler si suicida, resa senza condizioni \t governo provvisorio di Donitz \\
Nel 46, conferenza di parigi, germania firma pace \t prima c'è ialta dove si decide la divisione della germania in 4 zone e di berlino in 4 zone \\
Si decide anche denazificazione e quanto doveva pagare \t la conferenza di ialta conclude la ww2 ma non vede protagonista la germania \\
A ialta si decide anche di istituire un tribunale di guerra, per punire crimini contro l'umanita \t questo viene stabilito a ialta e ribadito a postdam \\

\sss[Processo di Norimberga]
Nell agosto del 45 urss, usa, uk e fr firmano accordo di londra, con cui viene istituito il tribunale militare internazionale che deve giudicare i nazisti \\
Ci sono 8 giudici civili, 4 giudici titolari e 4 transitori \t si sceglie Norimberga, perchè richiama le leggi di Norimberga \t si tiene volutamente qua, perche 10 anni prima erano state approvate le leggi razziali 10 anni prima \\
I capi di accusa possibili sono 4:
\bi
  \item cospirazione con fini aggressivi
  \item crimini contro la pace
  \item crimini di guerra
  \item crimini contro umanità
\ei
Vengono giudicati 22 nazisti, tra cui Goring e Hesserl \t capi + vicini \t qua viene anche pero processato Ribbentrop, che era un diplomatico \\
Il processo dura un anno, dal novembre del 45 al 46 \t ed è un tribunale equo: gli imputati possono difendersi, e molti dicono che non sono responsabili dei reati perchè hanno solo eseguito ordini dall'alto \\
1 ottobre del 46 pronunciata sentenza, che dichiara SS e tutte org. naziste org. criminali \\
12 di 22 vengono impiccati, 7 al carcere (tra 10 anni e ergastolo), e 3 assolti \t le condanne eseguite il 16 ottobre

\v

Questo processo è importante perche la violazione dei diritti umani viene sanzionata \t prima no \t vengono applicate le pene per i crimi contro l umanita x1 volta \\
Nel 48 vengono dichiarati i diritti fondamentali dall'onu

\v

In germania, ad ovest c'è zona occidentale che viene proclamata repubblica federale tedesca (riconosciuta ufficialmente nel 49) \\
Nel 47 le zone vengono unificate, e nel 49 si parla di rep. fed. ted \\
Ad est c'è repubblica democratica tedesca, stato comunista alleato dell'urss \t riconosciuto dall urss come l'unico stato tedesco legittimo, e non riconosce la repubblica federale \\
Anche quella democratica si costruisce nel 49 \\
Nel 46 si decide anche divisione di berlino, in est e ovest \t l'intera citta di berlino sta pero nella zona sovietica \\
Cosi si entra nella guerra fredda con berlino e germania divisa, con momenti di tensione \\
Nel 61 costruito il muro, e guerra fredda vede consolidamento anche fisico della separazione dei blocchi \\
Muro crolla nel novembre 89, e nel 1990 si tengono nella germania est le prime libere elezioni, e il 2 ottobre del 1990 si proclama l'unità della germania \\
La germania unita dara prova di grande solidita economica e politica \t est era molto + arretrata, e ovest si prende carico della situzione degradata dell est \t si fa carico del deficit disastroso dell'est

\ss[Austria]
Viene riconosciuta indipendente dopo la guerra, ma fino al 55 rimarrà divisa in 4 zone (sempre tra le 4 nazioni)

\ss[Giappone]
Occupato da usa fino al 51

\ss[Urss]
L'unione sovietica, alla fine della guerra, recupera diversi territori:
\bi
  \item ucraina
  \item bielorussia
  \item paesi baltici
  \item compensi territoriali da altri paesi, tra cui piccola parte della prussia orientale (dove c'è Koenigsberg, ora Kaliningrad)
\ei
Tra il 45 e il 48, urss costruisce un sistema comunista in tutta l europa orientale, che diventano stati satellite dell urss \\
Nel 47 nasce il COMINFORM (sciolto nel 56) = ufficio di informazione dei pariti comunisti, ovvero organismo che coordina l'azione dei partiti comunisti europei \\
Era indirizzato ai paesi satellite, ma aderiscono anche partiti dei paesi occidentali \\
L urss ottiene oltre vantaggi territoriali, anche vantaggi economici \t chiede ripagamento dei danni di guerra ai paesi con cui aveva combattuto \\
La germania dell'est era stata nemica, quindi ora deve pagare \t anche a ungheria, romania, bulgaria \\
Inoltre requisisce tutti i beni ex-tedeschi \t ovvero i beni dei cittadini tedeschi, quelli confiscati dai nazisti agli ebrei, etc. \\
Impone anche il pagamento del trasporto delle truppe sovietiche \t quando venivano mandate a presidiare regioni, queste regioni dovevano pagare il mantenimento e il trasporto di queste truppe \\
Il sistema economico grava pesantemente sui paesi satellite

\v

1947 c'è COMINFORM, mentre nel 49 c'è il COMECON = consiglio di assistenza economica, proposto contro il piano Marshall a cui aderiscono i paesi satellite \\
Si vuole creare grande alleanza economica \\
Nel 55 viene proposto il patto di Varsavia, che è alleanza militare contro NATO del 49 \\
Si arriva quindi a due blocchi contrapposti

\v

Alla fine arriva Gorbachov, che parla di Glasnost (trasparenza) e Perestroika (apretura) \t fa vacillare sistema comunista rigido \\
Trova interlocutore in Regan, anche se all'inizio è molto duro \t poi pero cambia linea \\
Gorbachov sara colui che portera l urss al suo disfacimento \t Gorbachov voleva ricostruire il sistema socialista, lo vuole ripensare dal punto di vista economico e nelle sue basi \t così non è vincente \\
Bisogna suscitare risposta nella popolazione, che era apatica però \t riceveva solo informazioni false dallo stato, e quindi non era spinta a partecipare attivamente \\
Quindi gorbachov vuole sucistare interesse nella politica, e dare notizie vere \\
Gorbachov incontra Regan, inizia dissoluzione urss \t nei paesi satellite si muovono \\
Crolla il muro di Berlini

\sss[Bulgaria] 
Mentre in Bulgaria nel 89 ci sono manifestazioni di protesta, nel 90 libere elezioni e nel 91 nuova costituzione

\sss[Ungheria] 
C'è Dubjeck, che guida proteste (insieme a Avel) \t il politburg da le dimissioni, sciopero contro il regime che porta al 28 dicembre dell ? con cui Avel diventa presidente repubblica e Dubjeck presidente del consiglio

\sss[Romania]
C'è Ciaogescu, che si era dichiarato nemico di Gorbachov \t ma non riesce a fermare il dissenso, nel dicembre 89 ci sono manifestazioni che pero vegono represse (in particolare a Temisoara) \\
Quasi da guerra civile tra popolazione e polizia \\
Poi comitato di salvezza processa in diretta TV Ciaogescu, e vengono fucilati sempre in diretta \\
Nel 90 questo comitato vince alle elezioni, ma continuano opposizioni (non piaceva molto questo comitato) \t e situazione rimane instabile

\sss[Jugoslavia]
Nel 90 si abbandona il monopolio del partito comunista \t e si cerca via democratica \\
Ci sono pero tensioni etniche nazionalistiche, che erano tenute sotto controllo dalla dittatura \t la democrazia invece ammorbidisce, e questo porta alla guerra e quindi la divisione della Jugoslavia

\sss[Albania]
Anche qua rivolte contro regime comunsta \t le prime elezioni libere sono nel 91

\v

Nel frattempo in urss cresce opposizione a Gorbachov \t Yeltsin guida i radicali che invece supportano Gorbachov \\
Yeltsin sara il tutore dell ordine e prendera posto di Gorbachov \\
Il partito, i servizi segreti e esercito tentano colpo di stato per tornare al comunismo pre Gorbachov \t lo tentato pero mentre Gorbachov per motivi di salute era in Crimea \\
Nell agosto del 91, il sedicente del comitato di emergenza (con grandi conservatori) \t il comitato dichiara che bisogna sostituire Gorbachov perche è troppo malato \\
A mosca pero c'era Yeltsin \t arrivano carri armati, ma Yeltsin è a capo del parlamento russo e con il suo supporto si oppone a questo tentativo di forza \\
I congiurati alla fine devono desistere: no appoggio popolazione e non avevano coesione \\
19 agosto c'è colpo di stato, ma poi vengono arrestati

\v

Gorbachov torna a mosca indebolito politicamente \\
Il 25 dicembre del 91 annuncera le sue dimissioni, e sancisce la fine dell'urss \t questa è la fine di un processo pero:
\bi
  \item settembre: si scioglie il congresso dei deputati del popolo
  \item 6 novembre: Yeltsin scoglie partito comunista russo e quello sovitico
  \item il 21 dicembre: i presidente di 11 delle 15 repubblica sovietiche, avevano firmato ad Almata la nascita della CSI (comunita degli stati indipendenti)
\ei
Nella CSI non entrano Georgia e neanche paesi balitci \t alla fine del 91 non piu urss, e questo cambia ordine internaizonale \t finisce anche il bipolarismo, che pero aveva garantito un ordine mondiale \\
Adesso invece c'è disordine mondiale \t ognuno per la sua strada \\
Alla fine degli anni 90, Yeltsin lascia il posto a Putin \t Yeltsin aveva problemi di salute e coinvolto in questioni giudiziari
\end{document}
